% Taller 1 · Filosofía 10.° · Periodo IV · «Mi quaestio» (semana 4, clase 2 de la parcelación)
% Imagen: img/tomas.jpg (Crivelli, 1476; Wikimedia Commons, dominio público).
% Desempeño: D3 (FIL-10-P4-D3). Una hora (≈55 min), en parejas, se resuelve en la hoja.
% Compilar: bash curriculo/plantilla/compilar.sh curriculo/filosofia/10/taller-1-P4.tex
\documentclass[taller]{planeacion}
\materia{filosofia}
\grado{10.°} \periodo{IV} \anio{2027}
\begin{document}
\fichataller{1}{Semana 4, clase 2}{55 min}{En parejas}

\titulotaller{Mi \emph{quaestio}}{En parejas. Todo se escribe \textbf{en esta hoja}; al final, entréguenla al docente
  para el sello.}

\reto{0}{El modelo: un artículo de la \emph{Suma teológica}}{10 min}{0,5}
\imagenlado[0.2]{img/tomas.jpg}{Tomás de Aquino (1225–1274) enseñó en la Universidad de París. Escribe cada
  artículo de la \emph{Suma teológica} como una disputa en la universidad: primero las objeciones,
  luego la autoridad contraria, su respuesta y la respuesta a cada objeción. Así lo hace
  (fragmentos):\par\vspace{3pt}
  {\footnotesize \emph{Santo Tomás de Aquino}, Carlo Crivelli (1476), National Gallery de Londres.
  Wikimedia Commons, dominio público.}}
\fuente{\textbf{¿Existe Dios?} \textbf{Objeción.} Parece que Dios no existe. Si uno de dos
  contrarios fuera infinito, destruiría por completo al otro. Con el nombre de Dios se entiende un
  bien infinito; si Dios existiera, no habría ningún mal. Pero en el mundo hay mal. Luego Dios no
  existe. \textbf{En sentido contrario}, dice el libro del Éxodo (3,~14): «Yo soy el que soy».
  \textbf{Respondo:} la existencia de Dios puede probarse por cinco vías; la primera es la del
  movimiento: todo lo que se mueve es movido por otro… \textbf{A la objeción:} como dice Agustín,
  Dios, sumamente bueno, no permitiría ningún mal en sus obras si no fuera tan poderoso y bueno que
  pudiera sacar bien del mismo mal.}
  {Tomás de Aquino, \emph{Suma teológica}, I, q. 2, a. 3 (fragmentos adaptados)}
\begin{tareas}
  \tarea{\envozalta lean el fragmento con su pareja.}
  \tarea{\enlahoja subrayen la objeción y encierren en un círculo la respuesta de Tomás a esa
    objeción. Fíjense: Tomás \textbf{empieza dándole la palabra a quien piensa distinto}.}
\end{tareas}

\reto{0}{Su pregunta}{5 min}{0,5}
\begin{tareas}
  \tarea{Escojan \textbf{una} pregunta: ¿Es lícito usar inteligencia artificial en las tareas?
    · ¿Debería prohibirse el celular en el colegio? · ¿Es justo exigir uniforme? · O una propia,
    aprobada por el docente.}
  \tarea{\enlahoja escríbanla en la línea «Pregunta» del punto 3.}
\end{tareas}

\reto{0}{Nuestra \emph{quaestio}}{25 min}{2,5}
\begin{tareas}
  \tarea{\enlahoja completen cada parte en orden, como Tomás: primero las objeciones, después su
    respuesta.}
\end{tareas}
\textbf{Pregunta:}\lineas{1}
\textbf{Objeción 1.} Parece que…\lineas{4}
\textbf{Objeción 2.} Además…\lineas{3}
\textbf{En sentido contrario} (una autoridad o una razón fuerte a favor de la otra postura):\lineas{3}
\textbf{Respondemos que…} (nuestra postura con sus razones):\lineas{6}
\textbf{A la objeción 1:}\lineas{3}
\textbf{A la objeción 2:}\lineas{3}

\reto{0}{La objeción de otra pareja}{10 min}{1,0}
\begin{tareas}
  \tarea{Intercambien la hoja con otra pareja: ella escribe \textbf{una objeción más} en las líneas de
    abajo.}
  \tarea{Recuperen su hoja y, \textbf{en las mismas líneas}, respondan esa objeción.}
\end{tareas}
\lineas{3}

\reto{0}{Cierre}{5 min}{0,5}
\begin{tareas}
  \tarea{\envozalta respondan: ¿por qué Tomás pone primero las objeciones?}
\end{tareas}

\apartado{Cómo se valora}
\begin{rubrica}
  \criterio{Estructura}{Las cinco partes completas y en orden}{Cinco partes, alguna incompleta}{Faltan una o dos partes}{No sigue la estructura}
  \criterio{Objeciones}{Fuertes: el contrario las firmaría}{Razonables}{Débiles o repetidas}{Ausentes o de paja}
  \criterio{Respuesta}{Postura clara con razones propias}{Postura con alguna razón}{Postura sin razones}{Sin postura}
  \criterio{Respuesta a objeciones}{Responde cada una con argumentos}{Responde con algún argumento}{Repite la postura}{No responde}
\end{rubrica}
\cierretaller{5,0}
\end{document}
